\documentclass[11pt]{article}
\usepackage[margin=1in]{geometry}
\usepackage{amsmath,amssymb,amsthm,mathtools}
\usepackage[hidelinks]{hyperref}
\newtheorem{proposition}{Proposition}
\newcommand{\R}{\mathbb R}
\newcommand{\C}{\mathbb C}
\newcommand{\sgn}{\operatorname{sgn}}
\title{A literature-implied answer to the signed Fourier RKBS conjecture}
\author{Run \texttt{fa\_banach\_001} --- lane 12}
\date{11 August 2026}

\begin{document}
\maketitle

\noindent\textbf{Status.}
\emph{Literature-implied answer: full existential conjecture, with a stronger
Wiener-class family}.  This is an agent-identified application of a known
feature-map theorem, not a claim of a new literature solution.

\section{Source question}

Xu and Ye's \emph{Generalized Mercer Kernels and Reproducing Kernel Banach
Spaces} (arXiv:1412.8663) says in Remark 4.6, printed page 69 (PDF page 75):
\begin{quote}
``So, we conjecture that a non-positive definite function could be used to
construct a RKBS by the Fourier transform.''
\end{quote}

The decisive supporting result is Lin--Zhang--Zhang,
arXiv:1901.01002, Theorem 2.3 (PDF pages 5--6).  In abbreviated form, it says
that Banach spaces $W_1,W_2$ with a continuous bilinear pairing and feature
maps $\Phi_j$ whose spans separate the opposite space produce a pair of
RKBSs, with reproducing kernel
\[
 K(x,y)=\langle\Phi_1(x),\Phi_2(y)\rangle_{W_1\times W_2}.
\]

\section{Direct identification}

We use the Fourier convention
\[
 \check h(t)=\int_{\R^d}h(\xi)e^{2\pi i t\cdot\xi}\,d\xi.
\]

\begin{proposition}
Let $h\in L^1(\R^d)$ be nonzero, real, and even, and set
$\phi=\check h$.  For every conjugate pair $1\le p,q\le\infty$, the
translation-invariant kernel
\[
 K(x,y)=\phi(x-y)
\]
is the reproducing kernel of a pair of (complex) RKBSs isometric to
$L^p(E)$ and $L^q(E)$, where $E=\{\xi:h(\xi)\ne0\}$.  It also has a real
form.  In particular, sign-changing $h$ give non-positive-definite examples.
\end{proposition}

\begin{proof}
On $E$, define
\[
 a=|h|^{1/q},\qquad b=\sgn(h)|h|^{1/p},
\]
with the usual endpoint convention $r^0=1$ on $E$.  Take
$W_1=L^q(E)$, $W_2=L^p(E)$ and the bilinear pairing
$\langle u,v\rangle=\int_Euv$.  The feature maps
\[
 \Phi_1(x)(\xi)=a(\xi)e^{2\pi i x\cdot\xi},\qquad
 \Phi_2(y)(\xi)=b(\xi)e^{-2\pi i y\cdot\xi}
\]
belong to $W_1,W_2$, respectively, and H\"older's inequality gives a
continuous pairing.  Since $ab=h$,
\[
 \langle\Phi_1(x),\Phi_2(y)\rangle
 =\int_Eh(\xi)e^{2\pi i(x-y)\cdot\xi}\,d\xi
 =\phi(x-y).
\]

It remains only to check the separating-density hypothesis of Theorem 2.3.
If $v\in L^p(E)$ annihilates every $\Phi_1(x)$, then $av\in L^1(E)$ and
its Fourier transform vanishes everywhere.  Fourier uniqueness gives
$av=0$ almost everywhere, hence $v=0$ because $a>0$ on $E$.  The same
argument applied to $ub\in L^1(E)$ proves the other density condition.
Theorem 2.3 therefore supplies the two RKBSs and both reproducing identities.

For a real version, restrict $L^p(E)$ and $L^q(E)$ to the closed real
subspaces of Hermitian functions $v(-\xi)=\overline{v(\xi)}$.  Evenness of
$h$ makes both feature maps Hermitian, the pairing real, and the same
Fourier-uniqueness argument nondegenerate.  Thus the theorem applies over the
real scalar field as well.
\end{proof}

\section{Explicit non-positive-definite kernel}

For $d=1$, take
\[
 \phi(t)=e^{-\pi t^2}-\frac12e^{-\pi t^2/16},\qquad
 h(\xi)=e^{-\pi\xi^2}-2e^{-16\pi\xi^2}.
\]
Then $h\in L^1(\R)$ is real, even, nonzero, and sign-changing: $h(0)=-1$,
whereas $h(\xi)>0$ for sufficiently large $|\xi|$.  By the uniqueness part
of Bochner's theorem, a continuous positive-definite $\phi$ cannot have this
signed Fourier density.  Hence $\phi$ is not positive definite, while the
proposition makes $\phi(x-y)$ a two-sided reproducing kernel for every
$p$-model above.  This answers the source's existential conjecture.

\section{Provenance and search bounds}

The supporting authors do not mention Xu--Ye's conjecture; the relation is the
direct identification above.  The four run indexes and the parsed arXiv corpus
were searched for the arXiv id and signed/non-PD Fourier RKBS terminology.
Bounded arXiv web searches inspected the exact source, arXiv:1901.01002, and
the nearby signed-measure Krein-space paper arXiv:2006.00247.  No explicit
all-$p$ statement was found, but Theorem 2.3 already proves it after the
displayed factorization.  Literature provenance therefore controls the packet
classification.

\section*{References}

\begin{enumerate}
\item Y. Xu and Q. Ye, \emph{Generalized Mercer Kernels and Reproducing
Kernel Banach Spaces}, arXiv:1412.8663, Remark 4.6.
\item R. Lin, H. Zhang, and J. Zhang, \emph{On Reproducing Kernel Banach
Spaces: Generic Definitions and Unified Framework of Constructions},
arXiv:1901.01002, Theorem 2.3.
\item F. Liu, X. Huang, Y. Chen, and J. A. K. Suykens, \emph{Fast Learning
in Reproducing Kernel Krein Spaces via Signed Measures}, arXiv:2006.00247.
\end{enumerate}

\end{document}
